\chapter{Что такое искусственный интеллект}
\label{ch:intro}

В этой главе мы дадим неформальное представление о том, какие задачи
решает искусственный интеллект, и наметим математический «инструментарий»,
без которого современный ИИ не существовал бы.

\section{Краткая история идеи}

\textit{[Раздел-заготовка для авторского текста: от Алана Тьюринга и
Дартмутской конференции 1956 года до глубокого обучения и больших
языковых моделей. Здесь полезно вставить временну́ю шкалу с помощью
\texttt{tikz} и набор гиперссылок на оригинальные статьи.]}

\section{Чему мы научимся в этой книге}

В книге читателю встретятся три «слоя» материала:
\begin{enumerate}
  \item \textbf{Математические основы}: линейная алгебра, элементы
        анализа, теория вероятностей и численные методы оптимизации.
  \item \textbf{Алгоритмы машинного обучения}: от линейной регрессии
        до многослойных нейронных сетей.
  \item \textbf{Современный ИИ}: трансформеры, большие языковые модели,
        генеративные методы.
\end{enumerate}

В главе~\ref{ch:opt} мы подробно изучим один из старейших и до сих пор
важнейших численных методов — \emph{метод касательных Ньютона}, который
используется и для решения уравнений, и для обучения современных
нейросетей.
